% fluids_bucket: the free surface of water in a spinning bucket (Vol II Ch 2, fig:fluids-bucket)
% In the frame turning with the bucket each parcel feels g down and the pseudo-force
% per unit mass omega^2 r outward.  The free surface is perpendicular to their
% resultant: dz/dr = omega^2 r/g, so z = omega^2 r^2/(2g), a paraboloid.
% Drawn to the challenge's numbers: R = 10 cm, omega = 10 rad/s gives
% Delta z = omega^2 R^2/(2g) = 51 mm, so Delta z/R = 0.51 (R = 1.8 cm here,
% Delta z = 0.92 cm, z = 0.284 r^2 in cm).
% A paraboloid's mean height over the disc is half its rim height, so the level at
% rest (dashed) lies Delta z/2 above the centre of the curved surface.
% At r = 0.95 cm the slope is 2(0.284)(0.95) = 0.54; arrows g (0.9) and
% omega^2 r (0.54 x 0.9 = 0.49) have a resultant perpendicular to the surface there.
\begin{tikzpicture}[x=1cm,y=1cm,
    water/.style={fill=cPrim!12,draw=none},
    declare function={zs(\x)=1.5+0.284*\x*\x;}]
  \def\R{1.8}
  % water
  \fill[water] (-\R,0) -- plot[domain=-\R:\R,samples=50] (\x,{zs(\x)}) -- (\R,0) -- cycle;
  \draw[cPrim,construction] plot[domain=-\R:\R,samples=50] (\x,{zs(\x)});
  % level at rest
  \draw[guide] (-\R,1.96) -- (\R,1.96);
  \node[anchor=east,cSteel,font=\footnotesize] at (-\R-0.1,1.96) {level at rest};
  % bucket
  \draw[outline] (-\R,3.0) -- (-\R,0) -- (\R,0) -- (\R,3.0);
  % axis and sense of rotation
  \draw[rotation axis] (0,-0.3) -- (0,3.5);
  \draw[rotation,cFor] (-0.45,3.25) arc[start angle=200,end angle=520,x radius=0.48,y radius=0.13];
  \node[anchor=west,cFor] at (0.5,3.3) {$\omega$};
  % rise at the rim
  \draw[dimtick] (0,1.5) -- (\R+0.55,1.5);
  \draw[dimtick] (\R,{zs(\R)}) -- (\R+0.55,{zs(\R)});
  \draw[dim] (\R+0.4,1.5) -- (\R+0.4,{zs(\R)}) node[midway,anchor=west,cSteel] {$\Delta z$};
  \draw[dim] (0,-0.5) -- (\R,-0.5) node[midway,anchor=north,cSteel] {$R$};
  % effective gravity at one point of the surface
  \coordinate (Q) at (-0.95,{zs(-0.95)});
  \draw[accel] (Q) -- ++(0,-0.9) node[anchor=north,cPrim] {$g$};
  \draw[accel] (Q) -- ++(-0.49,0);
  \node[anchor=north east,cPrim,inner sep=1pt] at ($(Q)+(-0.5,-0.04)$) {$\omega^2r$};
  \draw[accel,resultant] (Q) -- ++(-0.49,-0.9);
  \node[dot] at (Q) {};
\end{tikzpicture}
